\begin{lemma}[The theorem of sums for the regularized form, with a regularization-free bound]
Let $E$ be a finite-dimensional real inner-product space equipped with its Borel
$\sigma$-algebra, let $N\ge1$ and write $V=\ell^2\bigl(\bigoplus_{j=1}^NE\bigr)$
for the $\ell^2$-product of $N$ copies of $E$, whose elements are written
$Y=(y_j)_{j=1}^N$ and whose norm satisfies $|Y|^2=\sum_j|y_j|^2$. For each $j$
let $O_j\subseteq E$ be open and let $g_j:E\to\R$ be continuous on $O_j$. Let
$\Psi:V\to\R$, let $\widehat Z\in V$ with $\widehat z_j\in O_j$ for every $j$,
assume that $\Psi$ is $C^2$ near $\widehat Z$ and that
$Y\mapsto\sum_jg_j(y_j)-\Psi(Y)$ has a local maximum at $\widehat Z$. Set
$A=D^2\Psi(\widehat Z)$ and $\widehat p_j=D_j\Psi(\widehat Z)$, let
$\bar A:V\to V$ be the self-adjoint operator with
$\langle\bar A\xi,\zeta\rangle=A(\xi,\zeta)$, and let $\eta>0$. Then there is a
constant $\Lambda\ge0$ such that for every $\tau\in(0,1]$ there are symmetric
continuous bilinear forms $X_j^\tau$ on $E$ $(1\le j\le N)$ with
\begin{enumerate}
\item[(T1)] $\|X_j^\tau\|\le\Lambda$ for every $j$, with $\Lambda$ independent
  of $\tau$;
\item[(T2)] $\bigl(\widehat p_j,X_j^\tau\bigr)
  \in\overline J^{2,+}g_j(\widehat z_j)$ for every $j$, where
  $\overline J^{2,+}$ is the closed superjet of \noderef{ClosedSuperjet};
\item[(T3)] writing $B_\tau=\bar A+\tau\,\mathrm{id}$,
  \[
    \sum_jX_j^\tau[\zeta_j]
     \le\langle B_\tau\zeta,\zeta\rangle+\eta\,|B_\tau\zeta|^2
     =A[\zeta]+\tau|\zeta|^2+\eta\,|B_\tau\zeta|^2\qquad(\zeta\in V),
  \]
  where $X[h]=X(h,h)$, $A[\zeta]=A(\zeta,\zeta)$ and $|B_\tau\zeta|$ is the
  norm of the linear functional $\xi\mapsto A(\zeta,\xi)+\tau\langle\zeta,\xi\rangle$.
\end{enumerate}
\end{lemma}
\begin{proof}
For $a\in E$ and $1\le j\le N$ let $\iota_ja\in V$ be the vector with $j$-th
component $a$ and all other components $0$, so $|\iota_ja|=|a|$. Put
\[
  \widehat g_j(w)=g_j(\widehat z_j+w)-g_j(\widehat z_j)
   -\langle\widehat p_j,w\rangle\qquad(w\in E),
\]
a function defined on all of $E$ with $\widehat g_j(0)=0$.

\emph{Step 0: the constant $\Lambda$, chosen before $\tau$.}
For $\tau\in(0,1]$ the operator $B_\tau=\bar A+\tau\,\mathrm{id}$ is
self-adjoint with $\|B_\tau\|\le\|\bar A\|+\tau\le\|\bar A\|+1$. Put
\[
  \varepsilon(\tau)=\min\Bigl(\frac{\eta}{2},\
    \frac1{2\|B_\tau\|+2}\Bigr)>0,
  \qquad
  \Lambda=\max\Bigl(\frac2\eta,\ 2\|\bar A\|+4\Bigr)+4\bigl(\|\bar A\|+1\bigr).
\]
Then $\Lambda\ge0$, and $\Lambda$ depends only on $\eta$ and $\|\bar A\|$, not on
$\tau$. We record two consequences of the choice of $\varepsilon(\tau)$, writing
$\varepsilon=\varepsilon(\tau)$ from now on:
\begin{equation}\label{eq:srft-eps}
  \varepsilon\|B_\tau\|\le\frac{\|B_\tau\|}{2\|B_\tau\|+2}\le\tfrac12,
  \qquad
  2\varepsilon\le\eta .
\end{equation}
Moreover $\varepsilon^{-1}=\max\bigl(2/\eta,\ 2\|B_\tau\|+2\bigr)
\le\max\bigl(2/\eta,\ 2\|\bar A\|+4\bigr)$, so that
\begin{equation}\label{eq:srft-Lambda}
  \varepsilon^{-1}+2\cdot2\|B_\tau\|
   \le\max\Bigl(\frac2\eta,\ 2\|\bar A\|+4\Bigr)+4\bigl(\|\bar A\|+1\bigr)
   =\Lambda .
\end{equation}

Now fix $\tau\in(0,1]$.

\emph{Step 1: the quadratic test.}
By \noderef{SumRuleQuadraticTest}, applied with this $\tau$, there are $a>0$,
$C\ge0$ and continuous functions $G_j:E\to\R$ with $G_j(w)\le C$ for all $w$,
with $G_j(w)=\widehat g_j(w)-|w|^4$ for $|w|\le a$, and with
\begin{equation}\label{eq:srft-G3}
  \sum_jG_j(w_j)-\tfrac12\bigl(A[W]+\tau|W|^2\bigr)\le-\sum_j|w_j|^4
  \qquad(W\ne0).
\end{equation}
Since $\langle B_\tau W,W\rangle=\langle\bar AW,W\rangle+\tau|W|^2
=A[W]+\tau|W|^2$, the inequality \eqref{eq:srft-G3} is exactly
$\sum_jG_j(w_j)-\tfrac12\langle B_\tau W,W\rangle\le-\sum_j|w_j|^4$ for
$W\ne0$. Taking $w=0$ in the identity $G_j(w)=\widehat g_j(w)-|w|^4$ (legitimate
since $|0|=0\le a$) gives $G_j(0)=\widehat g_j(0)-0=0$.

\emph{Step 2: the resolvent.}
By \eqref{eq:srft-eps} we have $\varepsilon>0$ and
$\varepsilon\|B_\tau\|\le\tfrac12$, so \noderef{SumRuleResolvent}, applied in
$V$ to the self-adjoint operator $B_\tau$ with this $\varepsilon$, supplies a
self-adjoint continuous linear operator $A_\varepsilon:V\to V$ with the
properties (R1)--(R6) listed there; in particular
\begin{equation}\label{eq:srft-R}
  \|A_\varepsilon\|\le2\|B_\tau\|,
  \qquad
  \langle A_\varepsilon\zeta,\zeta\rangle
   \le\langle B_\tau\zeta,\zeta\rangle+\eta|B_\tau\zeta|^2\quad(\zeta\in V),
\end{equation}
the second inequality being (R6) with the present $\eta$, which is admissible
because $2\varepsilon\le\eta$ by \eqref{eq:srft-eps}.

\emph{Step 3: the duality test.}
The data assembled so far --- the continuous $G_j$ with $G_j\le C$, $C\ge0$ and
$G_j(0)=0$; the self-adjoint $B_\tau$ and $A_\varepsilon$; the parameter
$\varepsilon>0$; the identities (R2), (R3) and the positivity (R4) from Step 2;
and the strict inequality \eqref{eq:srft-G3} in the form of Step 1 --- are
exactly the hypotheses of \noderef{SumRuleDualityTest}. That lemma gives the
separability identity $\mathcal G^\varepsilon(Z)=\sum_jG_j^\varepsilon(z_j)$ for
$\mathcal G(W)=\sum_jG_j(w_j)$, and shows that
$\Upsilon(Z)=\mathcal G^\varepsilon(Z)
-\tfrac12\langle A_\varepsilon Z,Z\rangle$ satisfies $\Upsilon(0)=0$,
$\Upsilon(Z)<0$ for $Z\ne0$, is $\lambda$-semiconvex on $V$ with
$\lambda=\varepsilon^{-1}+\|A_\varepsilon\|$ in the sense of
\noderef{SemiconvexOn}, and is continuous.

\emph{Step 4: jets at shrinking scales.}
Fix $k\ge1$ and apply \noderef{SumRuleJetFromJensen} with the above
$G_j,C,\varepsilon,A_\varepsilon,\Upsilon$ and with $R=\delta'=1/k$. It produces
$\widehat Z'^{(k)}\in V$, $\Pi^{(k)}\in V$ and symmetric continuous bilinear
forms $X_j^{(k)}$ on $E$ with
\begin{align}
  \bigl|\widehat Z'^{(k)}\bigr|<\tfrac1k,
  &\qquad
  \bigl|\Pi^{(k)}\bigr|\le\tfrac1k+\|A_\varepsilon\|\tfrac1k
   =\frac{1+\|A_\varepsilon\|}{k},\label{eq:srft-small}\\
  \bigl\|X_j^{(k)}\bigr\|&\le\varepsilon^{-1}+2\|A_\varepsilon\|,
    \label{eq:srft-norm}\\
  \bigl(\Pi_j^{(k)},X_j^{(k)}\bigr)
   &\in J^{2,+}G_j\bigl(\widehat z'^{(k)}_j
     +\varepsilon\Pi_j^{(k)}\bigr),\label{eq:srft-jet}\\
  \sum_jX_j^{(k)}[\zeta_j]
   &\le\langle A_\varepsilon\zeta,\zeta\rangle\qquad(\zeta\in V).
    \label{eq:srft-form}
\end{align}
By \eqref{eq:srft-norm}, \eqref{eq:srft-R} and \eqref{eq:srft-Lambda},
\begin{equation}\label{eq:srft-unif}
  \bigl\|X_j^{(k)}\bigr\|\le\varepsilon^{-1}+2\|A_\varepsilon\|
   \le\varepsilon^{-1}+4\|B_\tau\|\le\Lambda
  \qquad(k\ge1,\ 1\le j\le N),
\end{equation}
a bound independent of both $k$ and $\tau$.

\emph{Step 5: the limit $k\to\infty$.}
By \eqref{eq:srft-unif} and \noderef{BilinearFormBolzanoWeierstrass} there are
continuous bilinear forms $X_j^\tau$ on $E$ and a strictly increasing
$\varphi:\N\to\N$ with $X_j^{(\varphi(k))}\to X_j^\tau$ for every $j$.

Each $X_j^\tau$ is symmetric: evaluation at a fixed pair of vectors is a
continuous linear functional on the space of continuous bilinear forms, so for
fixed $a,b\in E$ the real sequences $X_j^{(\varphi(k))}(a,b)$ and
$X_j^{(\varphi(k))}(b,a)$ converge to $X_j^\tau(a,b)$ and $X_j^\tau(b,a)$
respectively; these two sequences are equal term by term because each
$X_j^{(\varphi(k))}$ is symmetric, so their limits agree. Since the norm is
continuous, \eqref{eq:srft-unif} likewise passes to the limit and gives (T1):
$\|X_j^\tau\|\le\Lambda$.

Fix $\zeta\in V$. Evaluation at $(\zeta_j,\zeta_j)$ is continuous, so
$\sum_jX_j^{(\varphi(k))}[\zeta_j]\to\sum_jX_j^\tau[\zeta_j]$; passing to the
limit in \eqref{eq:srft-form} along the subsequence gives
\begin{equation}\label{eq:srft-limitform}
  \sum_jX_j^\tau[\zeta_j]\le\langle A_\varepsilon\zeta,\zeta\rangle
  \qquad(\zeta\in V).
\end{equation}

Next, put $z_k=\widehat z'^{(\varphi(k))}_j+\varepsilon\Pi_j^{(\varphi(k))}$,
$p_k=\Pi_j^{(\varphi(k))}$ and $\Xi_k=X_j^{(\varphi(k))}$ for a fixed $j$. From
\eqref{eq:srft-small} and $\varphi(k)\ge k$ we get
$\bigl|\widehat z'^{(\varphi(k))}_j\bigr|
\le\bigl|\widehat Z'^{(\varphi(k))}\bigr|<1/\varphi(k)\le1/k$ and
$|p_k|\le\bigl|\Pi^{(\varphi(k))}\bigr|\le(1+\|A_\varepsilon\|)/k$, so
$p_k\to0$ and $z_k\to0$. By \eqref{eq:srft-jet},
$(p_k,\Xi_k)\in J^{2,+}G_j(z_k)$; by continuity of $G_j$, $G_j(z_k)\to G_j(0)$;
and $\Xi_k\to X_j^\tau$. These are exactly the four requirements of
\noderef{ClosedSuperjet}, so
\begin{equation}\label{eq:srft-closed}
  \bigl(0,X_j^\tau\bigr)\in\overline J^{2,+}G_j(0)\qquad(1\le j\le N).
\end{equation}

\emph{Step 6: back to $g_j$.}
By Step 1, $G_j(w)=\widehat g_j(w)-|w|^4$ whenever $|w|\le a$, so in particular
the two functions $G_j$ and $\widehat g_j-|\cdot|^4$ agree on the ball
$\{|w|<a\}$ with $a>0$. By \noderef{ClosedSuperjetLocal},
$\overline J^{2,+}G_j(0)=\overline J^{2,+}\bigl(\widehat g_j-|\cdot|^4\bigr)(0)$,
and by \noderef{ClosedSuperjetQuartic} the latter equals
$\overline J^{2,+}\widehat g_j(0)$. Hence \eqref{eq:srft-closed} gives
$\bigl(0,X_j^\tau\bigr)\in\overline J^{2,+}\widehat g_j(0)$. Applying
\noderef{ClosedSuperjetAffineShift} with $z=\widehat z_j$, $p=\widehat p_j$,
$q=0$ and $X=X_j^\tau$ --- so that the normalized function there is exactly
$\widehat g_j$ --- converts this into
\[
  \bigl(\widehat p_j,X_j^\tau\bigr)\in\overline J^{2,+}g_j(\widehat z_j),
\]
which is (T2).

\emph{Step 7: the matrix inequality for $B_\tau$.}
Combining \eqref{eq:srft-limitform} with the second inequality of
\eqref{eq:srft-R} gives, for every $\zeta\in V$,
\[
  \sum_jX_j^\tau[\zeta_j]\le\langle A_\varepsilon\zeta,\zeta\rangle
   \le\langle B_\tau\zeta,\zeta\rangle+\eta|B_\tau\zeta|^2
   =A[\zeta]+\tau|\zeta|^2+\eta|B_\tau\zeta|^2,
\]
which is (T3). As $\tau\in(0,1]$ was arbitrary and $\Lambda$ was fixed in
Step 0 before $\tau$, the proof is complete.
\end{proof}
